\chapter{Methodology}
\label{ch:methodology}

\section{Introduction}
\label{sec:meth-intro}

This chapter presents the complete engineering methodology behind NirmalNode: the overall system architecture; the hardware used at the sensing and purification layers; the software running on the microcontroller; the communication and cloud architecture that connects the physical node to the AI layer; the Adaptive AI pipeline, including the mathematical formulation of the incremental-learning models used; and the decision logic that converts a pollution forecast into a physical purifier action. Each subsystem is described with enough implementation detail (equations, pseudocode, and architecture diagrams) to be reproducible.

\section{System Overview}
\label{sec:meth-overview}

NirmalNode is organized into five functional layers, illustrated in Figure~\ref{fig:five-layer}: (1) Environmental Monitoring, where physical sensors sample the air at the hotspot; (2) IoT Data Acquisition, where the ESP32 reads, cleans, timestamps, and transmits sensor data; (3) Cloud Data Storage, where historical and streaming data are persisted; (4) Adaptive AI Prediction, where an incrementally-updated model forecasts pollution 1--2 hours ahead; and (5) Local Air Purification, where the physical cyclone--HEPA--activated-carbon unit is triggered in response to the prediction.

\begin{figure}[htbp]
\centering
\begin{tikzpicture}[
  layer/.style={rectangle, draw=black, thick, rounded corners, fill=blue!8, minimum width=9cm, minimum height=1cm, font=\small},
  arr/.style={thick,-{Stealth[length=2.5mm]}}
]
\node[layer] (l1) {Layer 1 -- Environmental Monitoring (Sensors)};
\node[layer, below=0.55cm of l1] (l2) {Layer 2 -- IoT Data Acquisition (ESP32)};
\node[layer, below=0.55cm of l2] (l3) {Layer 3 -- Cloud Data Storage};
\node[layer, below=0.55cm of l3] (l4) {Layer 4 -- Adaptive AI Prediction (Incremental Learning)};
\node[layer, below=0.55cm of l4] (l5) {Layer 5 -- Local Air Purification (Cyclone + HEPA + Carbon)};
\draw[arr] (l1) -- (l2);
\draw[arr] (l2) -- (l3);
\draw[arr] (l3) -- (l4);
\draw[arr] (l4) -- (l5);
\end{tikzpicture}
\caption{Five-layer system architecture of NirmalNode}
\label{fig:five-layer}
\end{figure}

\section{System Flow}
\label{sec:meth-system-flow}

At a fixed sampling interval $\Delta t$ (prototype value: $\Delta t = 60$~s), the ESP32 reads all connected sensors, applies noise filtering, appends a timestamp and location identifier, and transmits the resulting record to the cloud database over Wi-Fi. The cloud layer appends the record to the historical dataset and forwards it to the Adaptive AI module, which produces two outputs: (a) an incremental update to its internal model parameters using the newly arrived data point, and (b) a forecast of pollutant concentration at time $t+h$ (with $h \in [1,2]$ hours). The forecast is passed to a threshold-based Decision Engine, which activates the purifier fan if the forecast crosses a hazard threshold, allowing purification to begin \emph{before} the hazardous condition materializes. Figure~\ref{fig:workflow} shows the complete, end-to-end workflow.

\begin{figure}[htbp]
\centering
\begin{tikzpicture}[
  node distance=0.42cm,
  box/.style={rectangle, draw=black, thick, rounded corners, fill=green!8, minimum width=7.2cm, minimum height=0.85cm, font=\small},
  arr/.style={thick,-{Stealth[length=2.2mm]}}
]
\node[box] (s1) {Environmental Sensors};
\node[box, below=of s1] (s2) {ESP32 (Read, Filter, Timestamp)};
\node[box, below=of s2] (s3) {Cloud Database};
\node[box, below=of s3] (s4) {Adaptive AI (Incremental Model)};
\node[box, below=of s4] (s5) {Pollution Prediction ($t+h$)};
\node[box, below=of s5] (s6) {Decision Engine (Threshold Logic)};
\node[box, below=of s6] (s7) {Purifier Control (Fan ON/OFF)};
\node[box, below=of s7] (s8) {Cyclone Separator};
\node[box, below=of s8] (s9) {HEPA Filter};
\node[box, below=of s9] (s10) {Activated Carbon Filter};
\node[box, below=of s10, fill=orange!12] (s11) {Cleaner Air for Worker};
\foreach \a/\b in {s1/s2,s2/s3,s3/s4,s4/s5,s5/s6,s6/s7,s7/s8,s8/s9,s9/s10,s10/s11}
  \draw[arr] (\a) -- (\b);
\end{tikzpicture}
\caption{Complete end-to-end system workflow, from sensing to clean-air delivery}
\label{fig:workflow}
\end{figure}

\FloatBarrier

\section{Hardware Architecture}
\label{sec:meth-hardware}

\subsection{ESP32 Microcontroller (Central Controller)}
\label{sec:meth-esp32}

The ESP32 Dev Module is used as the central controller of each NirmalNode unit. Its responsibilities are: (i) polling all attached sensors at the sampling interval $\Delta t$; (ii) applying a noise-filtering stage (Section~\ref{sec:meth-filtering}) to each raw reading; (iii) generating a timestamp and attaching the fixed location identifier of the node; (iv) formatting the resulting record (e.g., as JSON) for transmission; (v) transmitting the record to the cloud layer over Wi-Fi; (vi) receiving the prediction/decision result back from the cloud (or computing it locally, in an edge-AI configuration); and (vii) driving the purifier fan relay accordingly. The ESP32 was selected over alternative microcontrollers because it offers built-in Wi-Fi, multiple analog-to-digital converter (ADC) channels (needed for the analog gas sensors), UART interfaces (needed for the PMS5003's digital serial output), sufficiently high processing speed for on-device filtering and, where required, lightweight inference, low unit cost, low power consumption suitable for battery operation, and strong community/library support -- all of which directly support the Green IoT design goal of Section~\ref{sec:lr-green}.

\subsection{PMS5003 -- Laser Particulate Matter Sensor}
The PMS5003 is a laser-scattering particulate sensor that reports mass concentrations of PM1.0, PM2.5, and PM10. Air is drawn through an internal sensing chamber by a small fan; particles suspended in the air scatter a laser beam, and an internal photodiode measures the intensity and pattern of scattered light, from which the sensor's onboard processor estimates particle count and size distribution, and in turn mass concentration. PM2.5 is the single most safety-critical channel reported by this sensor, since particles below 2.5~$\mu$m are small enough to penetrate deep into the alveolar region of the lungs, which is the physiological basis for the health outcomes reviewed in Section~\ref{sec:lr-occupational-health}.

\subsection{MQ135 -- General Air Quality / Gas Sensor}
The MQ135 is a metal-oxide semiconductor (MOS) gas sensor whose resistance changes in the presence of a broad range of gases, including VOCs, smoke, NH$_3$, NO$_x$, alcohol vapour, and benzene, and which can additionally provide an indirect, non-quantitative estimate of elevated CO$_2$. Because it is a broad-spectrum sensor rather than a gas-specific one, its output is used in NirmalNode as a general air-quality indicator feature rather than as a calibrated single-gas concentration channel.

\subsection{MQ136 -- Hydrogen Sulphide Sensor}
The MQ136 is a MOS sensor specialized for hydrogen sulphide (H$_2$S) detection. H$_2$S is of particular concern in chemical processing, tanning/leather, waste-treatment, and food-processing industries -- all common in the Bangladeshi industrial landscape -- and is acutely toxic at elevated concentrations, motivating its inclusion as a dedicated sensing channel rather than relying on the broad-spectrum MQ135 alone.

\subsection{MiCS-4514 -- Dual-Element Gas Sensor}
The MiCS-4514 integrates two independent MOS sensing elements on a single package: a \emph{RED} element sensitive to carbon monoxide, hydrogen, methane, ethanol, and ammonia, and an \emph{OX} element sensitive to nitrogen dioxide. Because CO and NO$_2$ are both primary combustion by-products expected at boiler-room and generator-room hotspots, and because the two elements can be sampled simultaneously and independently, this sensor allows NirmalNode to monitor multiple industrially relevant gases from a single compact component.

\subsection{MP135 -- Supplementary Air Quality Sensor}
The MP135 provides a further general-purpose air-quality indication and is used to complement the MQ135 channel. Using multiple, partially overlapping gas sensors improves the robustness of the overall gas-sensing subsystem, since individual MOS sensors differ in their sensitivity range, response time, and cross-sensitivity to interfering gases; combining several sensors' outputs as separate input features allows the downstream AI model (Section~\ref{sec:meth-ai-pipeline}) to implicitly learn to weight and combine them appropriately for the specific hotspot in which the unit is deployed.

\subsection{BME280 (Optional) -- Temperature, Humidity, and Pressure}
An optional BME280 sensor provides ambient temperature, relative humidity, and atmospheric pressure readings. These environmental variables are included as AI input features because pollutant dispersion, particulate suspension, and gas-sensor response are all known to be modulated by temperature and humidity, and because diurnal pressure/temperature cycles correlate with certain industrial activity patterns (e.g., process heat cycles).

\subsection{Power Supply}
The bench-scale prototype is powered by a laboratory DC power supply for ease of development and testing. The design target for field deployment is a Lithium-Polymer (Li-Po) battery, selected to make the unit portable and independent of a fixed mains connection, which is important for retrofitting hotspots in existing factory floors where wiring a new mains-powered appliance may be impractical or unsafe.

\begin{table}[htbp]
\centering
\caption{Hardware components used in the NirmalNode prototype}
\label{tab:hardware-components}
\small
\begin{tabular}{@{}p{3.4cm}p{5.2cm}p{5.2cm}@{}}
\toprule
\textbf{Component} & \textbf{Function} & \textbf{Measured / Controlled Quantity} \\
\midrule
ESP32 Dev Module & Central controller & Sensor polling, filtering, communication, purifier control \\
PMS5003 & Laser particulate sensor & PM1.0, PM2.5, PM10 \\
MQ135 & General gas sensor & VOC, smoke, NH$_3$, NO$_x$, alcohol, benzene (indicative CO$_2$) \\
MQ136 & Specialized gas sensor & H$_2$S \\
MiCS-4514 & Dual-element gas sensor & CO, H$_2$, CH$_4$, ethanol, NH$_3$ (RED); NO$_2$ (OX) \\
MP135 & Supplementary gas sensor & General air-quality indication \\
BME280 (optional) & Environmental sensor & Temperature, humidity, pressure \\
DC Fan & Air mover & Draws air through the purification chamber \\
Cyclone Separator & Coarse pre-filter & Removes large dust particles \\
HEPA Filter & Fine particulate filter & PM2.5, PM1.0, dust, smoke, pollen, microorganisms \\
Activated Carbon Filter & Gas/odour filter & VOC, odour, organic gases, chemical vapours \\
Bench DC Supply / Li-Po Battery & Power source & Prototype bench power / portable field power \\
\bottomrule
\end{tabular}
\end{table}

\FloatBarrier

\section{Purification Subsystem}
\label{sec:meth-purification}

\subsection{Cyclone Separator}
The cyclone separator is the first purification stage. Polluted air enters the separator tangentially, and the resulting swirling flow subjects suspended particles to a centrifugal force
\begin{equation}
F_c = \frac{m\,v_t^2}{r}
\label{eq:centrifugal}
\end{equation}
where $m$ is particle mass, $v_t$ is the tangential velocity of the swirling air, and $r$ is the radius of the swirl path. Larger, heavier particles experience a proportionally larger $F_c$ and are thrown toward the outer wall of the separator, from which they fall into a dust-collection chamber, while the (now coarse-particle-depleted) air stream continues to the HEPA stage. This pre-separation step reduces the particulate loading reaching the HEPA filter, extending its usable lifetime.

\subsection{HEPA Filter}
The HEPA filter stage removes fine particulate matter, including PM2.5, PM1.0, residual dust, smoke particles, pollen, and many microorganisms. Filtration performance is conventionally expressed as a removal efficiency
\begin{equation}
\eta = 1 - \frac{C_{\text{out}}}{C_{\text{in}}}
\label{eq:hepa-efficiency}
\end{equation}
where $C_{\text{in}}$ and $C_{\text{out}}$ are the pollutant mass concentrations upstream and downstream of the filter, respectively. A filter meeting the HEPA standard achieves $\eta \geq 0.9997$ (i.e., at least 99.97\% removal) for particles of the most penetrating size, approximately 0.3~$\mu$m in diameter.

\subsection{Activated Carbon Filter}
The activated carbon stage removes gaseous pollutants -- VOCs, odour-causing compounds, and other organic/chemical vapours -- that pass through the HEPA stage unaffected, via physical adsorption onto the very large internal pore surface area of activated carbon. Adsorption capacity as a function of gas-phase concentration at equilibrium is commonly described by the Langmuir isotherm,
\begin{equation}
q_e = \frac{q_{\max} K C_e}{1 + K C_e}
\label{eq:langmuir}
\end{equation}
where $q_e$ is the mass of pollutant adsorbed per unit mass of carbon at equilibrium, $C_e$ is the equilibrium gas-phase concentration, $q_{\max}$ is the maximum adsorption capacity of the carbon, and $K$ is the Langmuir adsorption constant. This relationship is used qualitatively in this thesis to motivate periodic replacement of the activated-carbon cartridge as its effective $q_{\max}$ is approached, rather than as a fitted quantitative model (fitting Equation~\ref{eq:langmuir} to the specific carbon medium used is left as a calibration task; see Chapter 4, Limitations).

\subsection{Purification Pipeline}
Figure~\ref{fig:purification-pipeline} shows the physical purification pipeline in isolation, corresponding to Layer 5 of Figure~\ref{fig:five-layer}.

\begin{figure}[htbp]
\centering
\begin{tikzpicture}[
  node distance=0.5cm,
  pbox/.style={rectangle, draw=black, thick, rounded corners, fill=purple!8, minimum width=6.5cm, minimum height=0.85cm, font=\small},
  arr/.style={thick,-{Stealth[length=2.2mm]}}
]
\node[pbox] (p1) {Polluted Air Intake (DC Fan)};
\node[pbox, below=of p1] (p2) {Cyclone Separator (coarse dust removal)};
\node[pbox, below=of p2] (p3) {HEPA Filter (PM2.5 / PM1.0 removal)};
\node[pbox, below=of p3] (p4) {Activated Carbon Filter (VOC / odour removal)};
\node[pbox, below=of p4, fill=orange!12] (p5) {Purified Air Exit to Worker Zone};
\foreach \a/\b in {p1/p2,p2/p3,p3/p4,p4/p5}
  \draw[arr] (\a) -- (\b);
\end{tikzpicture}
\caption{Local air purification pipeline}
\label{fig:purification-pipeline}
\end{figure}

\FloatBarrier

\section{Software Architecture and Sensor Filtering}
\label{sec:meth-software}

\subsection{Noise Filtering}
\label{sec:meth-filtering}
Raw MOS gas-sensor readings are known to be noisy and drift-prone. Each raw reading $x_{\text{raw}}[k]$ at sample index $k$ is smoothed on-device using an exponential moving average (EMA),
\begin{equation}
\hat{x}[k] = \alpha\, x_{\text{raw}}[k] + (1-\alpha)\, \hat{x}[k-1]
\label{eq:ema}
\end{equation}
where $\alpha \in (0,1]$ is a smoothing constant chosen to trade off responsiveness against noise suppression (a higher $\alpha$ tracks rapid pollution changes more closely; a lower $\alpha$ suppresses more sensor noise). The filtered value $\hat{x}[k]$, rather than the raw reading, is used for both cloud transmission and AI model input.

\subsection{ESP32 Firmware Logic}
Algorithm~\ref{alg:esp32-loop} summarizes the main firmware loop executed on the ESP32.

\begin{algorithm}[htbp]
\caption{ESP32 Main Sensing and Transmission Loop}
\label{alg:esp32-loop}
\begin{algorithmic}[1]
\State Initialize Wi-Fi connection, sensor interfaces, and EMA state $\hat{x} \leftarrow 0$
\Loop
    \State Read raw values from PMS5003, MQ135, MQ136, MiCS-4514, MP135, BME280
    \For{each sensor channel $x_{\text{raw}}$}
        \State $\hat{x} \leftarrow \alpha\, x_{\text{raw}} + (1-\alpha)\, \hat{x}$ \Comment{EMA filtering, Eq.~\ref{eq:ema}}
    \EndFor
    \State Generate timestamp $t$ and attach fixed location ID
    \State Assemble record $r = \{\hat{x}_{\text{PM1.0..PM10}}, \hat{x}_{\text{CO,NO}_2}, \hat{x}_{\text{VOC,NH}_3,\text{H}_2\text{S}}, \hat{x}_{T,H,P}, t, \text{locationID}\}$ \Comment{full feature set, Eq.~\ref{eq:feature-vector}}
    \State Transmit $r$ to Cloud Database over Wi-Fi
    \State Receive Decision Engine output $d \in \{\text{ON}, \text{OFF}\}$ from cloud (or compute locally at the edge)
    \State Set purifier relay state according to $d$
    \State Sleep for sampling interval $\Delta t$
\EndLoop
\end{algorithmic}
\end{algorithm}

\FloatBarrier

\section{Communication and Cloud Architecture}
\label{sec:meth-cloud}

\subsection{Communication Architecture}
Each NirmalNode unit communicates with the cloud layer over Wi-Fi using a lightweight, JSON-formatted payload sent either via HTTPS POST requests or an MQTT publish to a node-specific topic (e.g., \texttt{nirmalnode/<locationID>/telemetry}). MQTT is preferred where available because its publish/subscribe model and small message overhead further support the Green IoT objective of minimizing communication energy cost, and because it naturally supports the Decision Engine subscribing to, and acting on, the same telemetry stream in near real time.

\subsection{Cloud Data Architecture}
The cloud layer maintains a time-series database in which every record is keyed by (\texttt{location\_id}, \texttt{timestamp}) and stores the full feature vector described in Section~\ref{sec:meth-data-pipeline}. This store serves three purposes: (i) it is the historical dataset from which the initial AI model (Section~\ref{sec:meth-ai-pipeline}) is trained; (ii) it is the source of the continuous stream of new records used for incremental model updates; and (iii) it supports a monitoring dashboard for human oversight of both raw pollutant levels and model predictions. Figure~\ref{fig:cloud-arch} depicts this architecture.

\begin{figure}[htbp]
\centering
\resizebox{\textwidth}{!}{%
\begin{tikzpicture}[
  node distance=0.9cm and 0.9cm,
  cbox/.style={rectangle, draw=black, thick, rounded corners, fill=cyan!8, minimum width=3.4cm, minimum height=1cm, font=\small, align=center},
  arr/.style={thick,-{Stealth[length=2.2mm]}}
]
\node[cbox] (node1) {NirmalNode Unit\\(Location A)};
\node[cbox, below=of node1] (node2) {NirmalNode Unit\\(Location B)};
\node[cbox, right=1.6cm of node1, yshift=-0.55cm] (mqtt) {MQTT / HTTPS Broker};
\node[cbox, right=1.6cm of mqtt, fill=yellow!12] (db) {Cloud Time-Series\\Database};
\node[cbox, below=of db] (ai) {Adaptive AI Service\\(Incremental Model)};
\node[cbox, right=1.6cm of db] (dash) {Monitoring Dashboard};
\draw[arr] (node1) -- (mqtt);
\draw[arr] (node2) -- (mqtt);
\draw[arr] (mqtt) -- (db);
\draw[arr] (db) -- (ai);
\draw[arr] (ai) -- (db);
\draw[arr] (db) -- (dash);
\draw[arr] (ai.east) to[bend left=25] (node1.south east);
\end{tikzpicture}%
}
\caption{Edge--cloud communication and data architecture (illustrated for two deployed units)}
\label{fig:cloud-arch}
\end{figure}

\FloatBarrier

\section{Data Pipeline and Feature Definition}
\label{sec:meth-data-pipeline}

Table~\ref{tab:features} lists the complete AI input feature vector, and Table~\ref{tab:ai-output} lists the AI output variables. At each sampling instant $t$, the feature vector is
\begin{equation}
\small
\begin{aligned}
\mathbf{x}_t = \big[\, &\text{PM1.0}_t,\ \text{PM2.5}_t,\ \text{PM10}_t,\ \text{CO}_t,\ \text{NO}_{2,t},\ \text{VOC}_t,\ \text{NH}_{3,t},\ \text{H}_2\text{S}_t, \\
&T_t,\ H_t,\ P_t,\ \text{hour}_t,\ \text{day}_t,\ y_{t-1} \,\big]
\end{aligned}
\label{eq:feature-vector}
\end{equation}
where $y_{t-1}$ is the previous observed pollution level (e.g., PM2.5 or a composite index), included so that the model can exploit short-term temporal autocorrelation in pollution levels in addition to the instantaneous multi-sensor reading.

\begin{table}[htbp]
\centering
\caption{AI input feature set}
\label{tab:features}
\small
\begin{tabular}{@{}p{4.5cm}p{8cm}@{}}
\toprule
\textbf{Feature} & \textbf{Description} \\
\midrule
PM1.0, PM2.5, PM10 & Particulate matter mass concentration (from PMS5003) \\
CO, NO$_2$ & Combustion-related gas concentration (from MiCS-4514) \\
VOC, NH$_3$ & Volatile organic compounds and ammonia (from MQ135 / MiCS-4514) \\
H$_2$S & Hydrogen sulphide (from MQ136) \\
Temperature, Humidity, Pressure & Ambient environmental conditions (from BME280) \\
Hour of Day & Cyclical feature capturing diurnal industrial activity patterns \\
Day of Week & Captures weekly production-schedule patterns \\
Previous Pollution Level ($y_{t-1}$) & Short-term temporal autocorrelation feature \\
\bottomrule
\end{tabular}
\end{table}

\begin{table}[htbp]
\centering
\caption{AI output variables}
\label{tab:ai-output}
\small
\begin{tabular}{@{}p{4.5cm}p{8cm}@{}}
\toprule
\textbf{Output} & \textbf{Description} \\
\midrule
Predicted pollution level $\hat{y}_{t+h}$ & Forecast pollutant concentration $h \in [1,2]$ hours ahead \\
Predicted AQI category & Discretized air-quality-index class derived from $\hat{y}_{t+h}$ \\
Risk level & Low / Moderate / High / Hazardous classification \\
Purifier recommendation & Binary decision (Activate / Do not activate) passed to the Decision Engine \\
\bottomrule
\end{tabular}
\end{table}

\FloatBarrier

\section{Adaptive AI Pipeline}
\label{sec:meth-ai-pipeline}

\subsection{Initial (Offline) Training}
Prior to field deployment, an initial model $f(\cdot\,;\mathbf{w}_0)$ is trained in batch mode on a historical Bangladesh air-quality dataset, using standard supervised regression, to obtain a reasonable starting parameter vector $\mathbf{w}_0$. This offline stage exists purely to give the deployed unit a sensible initial state; it is not repeated after deployment.

\subsection{Incremental (Online) Learning}
After deployment, the model is updated using \emph{incremental learning}: each new labelled observation $(\mathbf{x}_t, y_t)$ arriving from the field updates the model parameters directly, without reprocessing the full historical dataset. This thesis considers the following candidate incremental estimators, implemented via the River and/or scikit-learn (partial-fit) libraries.

\textbf{SGD Regressor.} Parameters are updated by a single stochastic gradient step per new sample,
\begin{equation}
\mathbf{w}_{t+1} = \mathbf{w}_t - \eta \, \nabla_{\mathbf{w}} \mathcal{L}\big(f(\mathbf{x}_t;\mathbf{w}_t), y_t\big)
\label{eq:sgd}
\end{equation}
where $\eta$ is the learning rate and $\mathcal{L}$ is a chosen loss function (e.g., squared error).

\textbf{Passive-Aggressive Regressor.} Parameters are updated only when the current model's error on the new sample exceeds a tolerance $\varepsilon$, with update magnitude
\begin{equation}
\mathbf{w}_{t+1} = \mathbf{w}_t + \tau_t\, \text{sign}(y_t - \hat{y}_t)\, \mathbf{x}_t, \qquad
\tau_t = \frac{\max(0,\, |y_t - \hat{y}_t| - \varepsilon)}{\lVert \mathbf{x}_t \rVert^2}
\label{eq:pa}
\end{equation}
which makes the model "passive" (no update) when its prediction is already acceptable and "aggressive" (a larger update) when the error is large -- a property well suited to a pollution signal that is mostly stable but occasionally shifts sharply.

\textbf{Hoeffding Tree.} A Hoeffding Tree grows a decision tree incrementally, splitting a leaf node once the Hoeffding bound guarantees, with high probability $1-\delta$, that the best-observed splitting feature at that leaf is indeed the true best split, without needing to revisit previously seen samples:
\begin{equation}
\epsilon = \sqrt{\frac{R^2 \ln(1/\delta)}{2n}}
\label{eq:hoeffding}
\end{equation}
where $R$ is the range of the relevant statistic, $n$ is the number of samples observed at the leaf, and $\epsilon$ is the resulting confidence bound on the estimated information-gain difference between the best and second-best candidate splits.

\textbf{Adaptive Random Forest (ARF).} An ensemble of incrementally-updated Hoeffding Trees, each augmented with a drift detector; when a constituent tree's error rate increases beyond a drift-detection threshold, that tree is replaced with a newly initialized tree, allowing the ensemble as a whole to adapt to \emph{concept drift} -- a change in the underlying relationship between features and pollution level, such as occurs when a unit is relocated to a new industrial area.

\subsection{Area Adaptation via Incremental Learning}
When a NirmalNode unit is relocated -- for example, from an industrial cluster in Gazipur (where pollution peaks may align with garment-production shift hours) to one in Narayanganj (where dyeing-industry effluent and emissions dominate) -- the incremental estimators above continue to update from the new location's data stream without any manual retraining step. The model's parameters (or, in the ARF case, its ensemble composition) gradually shift to reflect the new site's pollution signature purely as a byproduct of the same online-update mechanism used for routine continuous improvement, which is precisely the property that Section~\ref{sec:lr-adaptive-ai} identified as absent from the reviewed literature.

\subsection{Prediction Process}
At each sampling instant $t$, the current model produces a forecast
\begin{equation}
\hat{y}_{t+h} = f(\mathbf{x}_t;\, \mathbf{w}_t)
\label{eq:forecast}
\end{equation}
for a look-ahead horizon $h \in [1,2]$ hours. This forecast, rather than the instantaneous reading $y_t$, is what is passed to the Decision Engine (Section~\ref{sec:meth-decision}), which is the mechanism by which NirmalNode achieves \emph{predictive}, rather than reactive, purifier activation.

\subsection{Evaluation Metrics}
The predictive performance of the incremental model is evaluated in Chapter 4 using:
\begin{align}
\text{MAE} &= \frac{1}{n}\sum_{i=1}^{n} \left| y_i - \hat{y}_i \right| \label{eq:mae}\\
\text{RMSE} &= \sqrt{\frac{1}{n}\sum_{i=1}^{n} \left( y_i - \hat{y}_i \right)^2} \label{eq:rmse}\\
R^2 &= 1 - \frac{\sum_{i=1}^{n}(y_i - \hat{y}_i)^2}{\sum_{i=1}^{n}(y_i - \bar{y})^2} \label{eq:r2}
\end{align}
and, for the derived risk-level classification task,
\begin{align}
\text{Precision} &= \frac{TP}{TP+FP}, \qquad
\text{Recall} = \frac{TP}{TP+FN}, \qquad
\text{F1} = \frac{2\cdot\text{Precision}\cdot\text{Recall}}{\text{Precision}+\text{Recall}}
\label{eq:prf1}
\end{align}
where $TP$, $FP$, and $FN$ denote true positives, false positives, and false negatives with respect to correctly forecasting a "hazardous" risk level.

\section{Purifier Decision Engine}
\label{sec:meth-decision}

The Decision Engine converts the AI's forecast $\hat{y}_{t+h}$ into a binary purifier command using a threshold rule. Let $\tau_{\text{hazard}}$ be the pollutant concentration threshold above which conditions are considered hazardous for a worker at the monitored hotspot (e.g., a PM2.5 or composite-index threshold derived from relevant air-quality guideline levels). The decision is
\begin{equation}
d_t =
\begin{cases}
\text{ON}, & \hat{y}_{t+h} \geq \tau_{\text{hazard}} \\
\text{OFF}, & \hat{y}_{t+h} < \tau_{\text{hazard}}
\end{cases}
\label{eq:decision-rule}
\end{equation}
Algorithm~\ref{alg:decision-engine} summarizes this logic together with the incremental model update that accompanies every new observation.

\begin{algorithm}[htbp]
\caption{Adaptive AI Update and Purifier Decision Engine}
\label{alg:decision-engine}
\begin{algorithmic}[1]
\Require Incoming record $r_t = (\mathbf{x}_t, y_t)$, current model $f(\cdot\,;\mathbf{w}_t)$, hazard threshold $\tau_{\text{hazard}}$, horizon $h$
\State $\hat{y}_{t+h} \leftarrow f(\mathbf{x}_t; \mathbf{w}_t)$ \Comment{Forecast using current model, Eq.~\ref{eq:forecast}}
\If{$\hat{y}_{t+h} \geq \tau_{\text{hazard}}$}
    \State $d_t \leftarrow \text{ON}$
\Else
    \State $d_t \leftarrow \text{OFF}$
\EndIf
\State Send $d_t$ to ESP32 purifier relay
\State $\mathbf{w}_{t+1} \leftarrow \text{IncrementalUpdate}(\mathbf{w}_t, \mathbf{x}_t, y_t)$ \Comment{Eq.~\ref{eq:sgd} or \ref{eq:pa}, or ARF/Hoeffding update}
\State Append $r_t$ to Cloud Database
\State \Return $d_t$, $\mathbf{w}_{t+1}$
\end{algorithmic}
\end{algorithm}

\section{Additional System Diagrams}
\label{sec:meth-extra-diagrams}

The preceding sections presented the overall system architecture (Figure~\ref{fig:five-layer}), end-to-end workflow (Figure~\ref{fig:workflow}), purification pipeline (Figure~\ref{fig:purification-pipeline}), and cloud/communication architecture (Figure~\ref{fig:cloud-arch}). This section presents the remaining architecture-level diagrams -- the hardware block diagram, the software block diagram, the AI pipeline, the incremental-learning pipeline, the purifier decision flowchart, and the multi-hotspot deployment architecture -- that together give a complete visual specification of NirmalNode.

\subsection{Hardware Block Diagram}
Figure~\ref{fig:hardware-block} shows the physical wiring-level block diagram of a single NirmalNode unit: each sensor connects to the ESP32 over its native interface (UART for the PMS5003, analog ADC channels for the MOS gas sensors, I$^2$C for the optional BME280), and the ESP32 drives the purifier fan through a relay/MOSFET switching stage.

\begin{figure}[htbp]
\centering
\resizebox{0.95\textwidth}{!}{%
\begin{tikzpicture}[
  node distance=0.5cm and 1.2cm,
  hbox/.style={rectangle, draw=black, thick, rounded corners, fill=blue!8, minimum width=3.4cm, minimum height=0.9cm, font=\small, align=center},
  ctrl/.style={rectangle, draw=black, thick, rounded corners, fill=orange!15, minimum width=3.6cm, minimum height=2.6cm, font=\bfseries\small, align=center},
  arr/.style={thick,-{Stealth[length=2.2mm]}}
]
\node[ctrl] (esp) {ESP32\\Dev Module\\(Central Controller)};
\node[hbox, above left=0.1cm and 1.4cm of esp] (pms) {PMS5003\\(UART)};
\node[hbox, left=1.4cm of esp] (mics) {MiCS-4514\\(Analog ADC)};
\node[hbox, below left=0.1cm and 1.4cm of esp] (mq) {MQ135 / MQ136\\(Analog ADC)};
\node[hbox, above=0.4cm of pms] (mp) {MP135\\(Analog ADC)};
\node[hbox, below=0.4cm of mq] (bme) {BME280\\(I$^2$C, optional)};
\node[hbox, right=1.4cm of esp, fill=red!10] (relay) {Relay / MOSFET\\Switching Stage};
\node[hbox, right=1.4cm of relay, fill=green!10] (fan) {12V DC Fan\\(Purifier Actuator)};
\node[hbox, above=0.4cm of relay, fill=gray!12] (wifi) {Wi-Fi Radio\\(to Cloud)};
\draw[arr] (pms) -- (esp);
\draw[arr] (mics) -- (esp);
\draw[arr] (mq) -- (esp);
\draw[arr] (mp) -- (esp);
\draw[arr] (bme) -- (esp);
\draw[arr] (esp) -- (relay);
\draw[arr] (relay) -- (fan);
\draw[arr] (esp) -- (wifi);
\end{tikzpicture}%
}
\caption{Hardware block diagram of a single NirmalNode unit}
\label{fig:hardware-block}
\end{figure}

\FloatBarrier

\subsection{Software Block Diagram}
Figure~\ref{fig:software-block} decomposes the system's software into its constituent modules, spanning the ESP32 firmware (device-side) and the cloud service (server-side).

\begin{figure}[htbp]
\centering
\resizebox{\textwidth}{!}{%
\begin{tikzpicture}[
  node distance=0.5cm and 0.7cm,
  sbox/.style={rectangle, draw=black, thick, rounded corners, fill=purple!8, minimum width=3.3cm, minimum height=1cm, font=\small, align=center},
  grp/.style={draw=black, dashed, inner sep=8pt, rounded corners},
  arr/.style={thick,-{Stealth[length=2.2mm]}}
]
\node[sbox] (drv) {Sensor Driver Layer};
\node[sbox, right=of drv] (filt) {EMA Noise Filter\\(Eq.~\ref{eq:ema})};
\node[sbox, right=of filt] (pack) {Record Packetizer\\(JSON, Eq.~\ref{eq:feature-vector})};
\node[sbox, right=of pack] (comm) {Communication Module\\(MQTT / HTTPS)};
\node[sbox, below=1.1cm of pack, fill=cyan!10] (store) {Cloud Storage Service};
\node[sbox, right=of store, fill=yellow!12] (aiserv) {Adaptive AI Service\\(Incremental Model)};
\node[sbox, right=of aiserv, fill=red!8] (dec) {Decision Engine\\(Algorithm~\ref{alg:decision-engine})};
\node[sbox, below=1.1cm of drv, fill=green!10] (act) {Actuator Control Module};
\draw[arr] (drv) -- (filt);
\draw[arr] (filt) -- (pack);
\draw[arr] (pack) -- (comm);
\draw[arr] (comm) -- (store);
\draw[arr] (store) -- (aiserv);
\draw[arr] (aiserv) -- (dec);
\draw[arr] (dec) to[bend right=15] (act);
\draw[arr] (act) -- (drv);
\node[grp, fit=(drv)(filt)(pack)(comm), label=above:{\small \textbf{ESP32 Firmware (Device Side)}}] {};
\node[grp, fit=(store)(aiserv)(dec), label=above:{\small \textbf{Cloud Service (Server Side)}}] {};
\end{tikzpicture}%
}
\caption{Software block diagram, spanning device-side firmware and server-side cloud service}
\label{fig:software-block}
\end{figure}

\FloatBarrier

\subsection{AI Pipeline Diagram}
Figure~\ref{fig:ai-pipeline} illustrates the complete AI pipeline, from historical pre-training through continuous incremental updating, described mathematically in Section~\ref{sec:meth-ai-pipeline}.

\begin{figure}[htbp]
\centering
\begin{tikzpicture}[
  node distance=0.5cm,
  aibox/.style={rectangle, draw=black, thick, rounded corners, fill=yellow!10, minimum width=8cm, minimum height=0.85cm, font=\small},
  arr/.style={thick,-{Stealth[length=2.2mm]}}
]
\node[aibox] (a1) {Historical Bangladesh Air-Quality Dataset};
\node[aibox, below=of a1] (a2) {Initial Batch Training (Baseline Model $\mathbf{w}_0$)};
\node[aibox, below=of a2, fill=orange!12] (a3) {Deployed Incremental Model $f(\cdot\,;\mathbf{w}_t)$};
\node[aibox, below=of a3] (a4) {Forecast $\hat{y}_{t+h}$ (Eq.~\ref{eq:forecast})};
\node[aibox, below=of a4] (a5) {Decision Engine (Section~\ref{sec:meth-decision})};
\node[aibox, below=of a5, fill=cyan!10] (a6) {Incremental Model Update (Eq.~\ref{eq:sgd} / \ref{eq:pa})};
\draw[arr] (a1) -- (a2);
\draw[arr] (a2) -- (a3);
\draw[arr] (a3) -- (a4);
\draw[arr] (a4) -- (a5);
\draw[arr] (a5) -- (a6);
\draw[arr] (a6.east) to[bend left=45] (a3.east);
\end{tikzpicture}
\caption{AI pipeline: from historical pre-training to continuous incremental updating}
\label{fig:ai-pipeline}
\end{figure}

\FloatBarrier

\subsection{Incremental Learning Pipeline}
Figure~\ref{fig:incremental-pipeline} details the prequential (test-then-train) loop that gives NirmalNode its Adaptive AI and area-specific-learning properties, as introduced conceptually in Section~\ref{sec:lr-adaptive-ai}.

\begin{figure}[htbp]
\centering
\begin{tikzpicture}[
  node distance=0.9cm and 0.9cm,
  ilbox/.style={rectangle, draw=black, thick, rounded corners, fill=green!8, minimum width=3.6cm, minimum height=1cm, font=\small, align=center},
  dcbox/.style={diamond, draw=black, thick, fill=red!8, minimum width=2.6cm, minimum height=1.2cm, font=\small, align=center, inner sep=1pt},
  arr/.style={thick,-{Stealth[length=2.2mm]}}
]
\node[ilbox] (i1) {New Sensor Record $(\mathbf{x}_t, y_t)$ Arrives};
\node[ilbox, below=of i1] (i2) {Predict $\hat{y}_{t+h}$ Using Current Model $\mathbf{w}_t$};
\node[ilbox, below=of i2] (i3) {Score Prediction Once True $y_{t+h}$ is Observed};
\node[ilbox, below=of i3] (i4) {Incrementally Update Model: $\mathbf{w}_t \rightarrow \mathbf{w}_{t+1}$};
\node[dcbox, right=1.6cm of i2] (drift) {Relocated / Drift Detected?};
\node[ilbox, right=1.6cm of drift, fill=orange!12] (adapt) {Area-Specific Re-Adaptation};
\draw[arr] (i1) -- (i2);
\draw[arr] (i2) -- (i3);
\draw[arr] (i3) -- (i4);
\draw[arr] (i4.west) to[bend right=60] (i1.west);
\draw[arr] (i2) -- (drift);
\draw[arr] (drift) -- node[above,font=\scriptsize]{Yes} (adapt);
\draw[arr] (adapt.south) to[bend left=40] (i4.east);
\draw[arr] (drift.south) to[bend right=25] node[right,font=\scriptsize]{No} (i3.east);
\end{tikzpicture}
\caption{Incremental (online) learning pipeline, including area-specific re-adaptation on relocation}
\label{fig:incremental-pipeline}
\end{figure}

\FloatBarrier

\subsection{Purifier Decision Flowchart}
Figure~\ref{fig:decision-flowchart} presents the threshold-based purifier decision logic of Section~\ref{sec:meth-decision} (Equation~\ref{eq:decision-rule}, Algorithm~\ref{alg:decision-engine}) as a visual flowchart.

\begin{figure}[htbp]
\centering
\begin{tikzpicture}[
  node distance=0.7cm,
  dfbox/.style={rectangle, draw=black, thick, rounded corners, fill=blue!8, minimum width=6.5cm, minimum height=0.85cm, font=\small},
  dfdec/.style={diamond, draw=black, thick, fill=red!10, minimum width=3.6cm, minimum height=1.4cm, font=\small, align=center, inner sep=1pt},
  arr/.style={thick,-{Stealth[length=2.2mm]}}
]
\node[dfbox] (f1) {Obtain Forecast $\hat{y}_{t+h}$};
\node[dfdec, below=of f1] (f2) {$\hat{y}_{t+h} \geq \tau_{\text{hazard}}$?};
\node[dfbox, below left=1cm and -1.2cm of f2, fill=green!12] (f3) {Purifier OFF};
\node[dfbox, below right=1cm and -1.2cm of f2, fill=orange!15] (f4) {Purifier ON (Pre-emptive Activation)};
\node[dfbox, below=1.4cm of f2, fill=cyan!10] (f5) {Log Decision + Update Model (Eq.~\ref{eq:sgd}/\ref{eq:pa})};
\draw[arr] (f1) -- (f2);
\draw[arr] (f2) -- node[above,font=\scriptsize]{No} (f3);
\draw[arr] (f2) -- node[above,font=\scriptsize]{Yes} (f4);
\draw[arr] (f3) |- (f5);
\draw[arr] (f4) |- (f5);
\end{tikzpicture}
\caption{Purifier decision flowchart}
\label{fig:decision-flowchart}
\end{figure}

\FloatBarrier

\subsection{Deployment Architecture}
Figure~\ref{fig:deployment-arch} shows how multiple NirmalNode units are envisaged to be deployed across the distinct hotspots of a single factory, each reporting independently to the shared cloud layer of Figure~\ref{fig:cloud-arch}, consistent with the multi-node future-work direction outlined in Section~\ref{sec:future-smartcity}.

\begin{figure}[htbp]
\centering
\begin{tikzpicture}[
  node distance=0.6cm and 1cm,
  dbox/.style={rectangle, draw=black, thick, rounded corners, fill=blue!8, minimum width=3.2cm, minimum height=0.8cm, font=\scriptsize, align=center},
  factory/.style={draw=black, thick, dashed, rounded corners, inner sep=10pt},
  arr/.style={thick,-{Stealth[length=2mm]}}
]
\node[dbox] (h1) {NirmalNode\\Welding Bay};
\node[dbox, right=of h1] (h2) {NirmalNode\\Boiler Room};
\node[dbox, right=of h2] (h3) {NirmalNode\\Generator Room};
\node[dbox, below=1.4cm of h2] (h4) {NirmalNode\\Chemical Store};
\node[dbox, below left=1.4cm and 0.2cm of h1] (h5) {NirmalNode\\Factory Entrance};
\node[dbox, below right=1.4cm and 0.2cm of h3] (h6) {NirmalNode\\Worker Operating Zone};
\node[factory, fit=(h1)(h2)(h3)(h4)(h5)(h6), label=above:{\small \textbf{Factory Floor -- Multiple Hotspots}}] (fbox) {};
\node[dbox, below=1.8cm of h4, fill=yellow!15, minimum width=4.5cm] (cloud) {Shared Cloud Platform\\(Storage, Adaptive AI, Dashboard)};
\draw[arr] (h4.south) -- (cloud.north);
\draw[arr] (h5.south) to[bend right=15] (cloud.west);
\draw[arr] (h6.south) to[bend left=15] (cloud.east);
\end{tikzpicture}
\caption{Deployment architecture: multiple NirmalNode units across a factory's pollution hotspots, sharing a common cloud platform}
\label{fig:deployment-arch}
\end{figure}

\FloatBarrier

\section{Chapter Summary}
\label{sec:meth-summary}

This chapter has presented the full NirmalNode methodology: a five-layer system architecture; hardware comprising an ESP32 controller, a PMS5003 particulate sensor, MQ135/MQ136/MiCS-4514/MP135 gas sensors, an optional BME280 environmental sensor, and a cyclone--HEPA--activated-carbon purification train; an EMA-filtered sensing firmware loop; an MQTT/HTTPS-based edge--cloud communication and storage architecture; an Adaptive AI pipeline built on incremental-learning estimators (SGD Regressor, Passive-Aggressive Regressor, Hoeffding Tree, Adaptive Random Forest) that update continuously and adapt to new deployment areas without full retraining; and a threshold-based Decision Engine that converts a 1--2 hour-ahead pollution forecast into a predictive, rather than reactive, purifier activation command. Chapter 4 now describes the experimental setup used to evaluate this methodology and presents the corresponding results.
